\documentclass[12pt]{article}
%\renewcommand\normalsize{\fontsize{10pt}{12pt}\selectfont}
\usepackage{graphicx} % Required for inserting images
\usepackage[left=2cm,right=2cm,top=2cm,bottom=1.5cm]{geometry}
\pagestyle{empty}

\title{Quiz 3}
%\author{Dennis Wilkinson}
%\date{March 31 2025; due Monday April 7}

\begin{document}

\setlength{\parindent}{0pt}
\setlength{\parskip}{3mm}

FS1 Wilkinson Quiz 3, April 24, 2025

1. Mars' fastest speed in its orbit occurs at a point we'll call A, and its slowest speed at a point B. 

a. At which point, A or B, is the orbital (i) kinetic energy, (ii) potential energy, and (iii) radius greater? Briefly justify your answer. (Here we define ``greater" PE to mean least negative, so a PE of -5 would be greater than a PE of -10.)


\bigskip
\bigskip
\bigskip
\bigskip
\bigskip
\bigskip
\bigskip
\bigskip
\bigskip


b. If you knew the orbital radius at point A, and the speeds at A and B, explain how you would calculate the orbital radius at point B. Words or equations okay.

\bigskip
\bigskip
\bigskip
\bigskip
\bigskip
\bigskip
\bigskip
\bigskip
\bigskip


2. In the future, perhaps we will launch spaceships from stations which are in orbit around Earth. 

a. In terms of the ship's energy, explain how launching from such a station would reduce the escape velocity needed to get away from Earth's gravity, as compared to launching an identical ship from Earth's surface.


\bigskip
\bigskip
\bigskip
\bigskip
\bigskip
\bigskip
\bigskip
\bigskip
\bigskip

b. If the station is located a distance $R/2$ above the Earth, where $R$ is the Earth's radius, find the escape velocity (to escape Earth) for a ship blasting off from the station. Escape $v$ from Earth is 11.2 km/s. 

\end{document}